\section{Results}


\subsection*{Physiologically grounded synthetic signals provide controlled ground truth for fidelity evaluation}

Evaluating digital twin fidelity requires signals whose ground-truth amplitude, frequency, and phase are fully specified, yet whose dynamics reflect the complexity of real biosignals \cite{clifford2006advanced, sornmo2005bioelectrical}. To achieve this, we fit closed-form parametric models to real physiological recordings from three clinical-grade datasets and sampled new signals from the resulting empirical parameter distributions (Methods; Supplementary Tables~S1--S2). Each signal family is defined by a closed-form parametric equation designed to isolate specific dynamical properties encountered in clinical biosignals. The single phase-modulated (SPM) harmonic family generates a sinusoidal carrier whose instantaneous phase is modulated by a second, slower sinusoid. This structure was chosen because it reproduces the way a single physiological rhythm, such as the cardiac cycle, varies in timing under the influence of a slower regulatory process like respiration, producing the phase variability characteristic of basic heart rate variability. The dual phase-modulated (DPM) harmonic family sums two independently modulated carrier-modulator pairs, each with its own frequency and modulation depth. This extension is necessary because neural and complex cardiac signals contain multiple co-occurring rhythmic components, such as simultaneous alpha and theta oscillations in EEG or overlapping atrial and ventricular dynamics in ECG, whose independent phase variations create the spectral complexity that simpler single-modulator models cannot reproduce. The drift-harmonic (DH) family generates a sinusoid whose amplitude decays exponentially over time and whose baseline shifts through a linear trend. This structure captures the slow wander and amplitude drift observed in peripheral cardiovascular signals such as photoplethysmograms, where sensor coupling changes and vasomotor fluctuations cause the signal envelope to evolve independently of the underlying oscillation. Full parametric equations are provided in Methods; parameter bounds derived from physiological fitting are reported in Supplementary Table~S1.

The drift-harmonic family, derived from blood volume pulse segments in the PPG-DaLiA dataset \cite{ppg}, captured the quasi-sinusoidal morphology of peripheral cardiovascular signals including slow amplitude drift and phase variation (Fig.~\ref{fig:fitting}C). Fitted frequency estimates concentrated near 1.0--1.1~Hz, consistent with resting heart rate, and drift coefficients remained near zero, confirming physiological plausibility (Supplementary Table~S2). The single and dual phase-modulated harmonic families, derived from the MIT-BIH Arrhythmia Database \cite{moody2001impact, goldberger2000physiobank}, captured the smooth oscillatory baseline of the S-Q interval, the segment between successive QRS complexes where atrial repolarization, the T~wave, and diastolic baseline dominate, while excluding sharp R-peak transients that violate sinusoidal decomposition assumptions (Fig.~\ref{fig:fitting}B). Carrier frequencies clustered near 1.0~Hz across all 48 records, with modulation depths and frequencies spanning ranges consistent with heart rate variability. The dual phase-modulated family, fitted to channel FP1-F7 from the CHB-MIT Scalp EEG Database \cite{shoeb2009application}, captured dominant cortical rhythmic activity while accommodating transient deflections characteristic of EEG recordings (Fig.~\ref{fig:fitting}A).

Because every synthetic signal is generated from closed-form equations with fully specified parameters, the ground-truth amplitude, instantaneous frequency, and instantaneous phase are analytically known at each time point. The generation procedure, including parameter uniqueness guarantees and dataset split sizes, is detailed in Supplementary Section~S3. Table~\ref{tab:table1} maps each synthetic paradigm to the physiological phenomenon it represents, establishing the clinical grounding for the evaluation paradigms used throughout this study.

\begin{figure}
    \centering
    \includegraphics[width=0.9\linewidth]{Figures/Framework/Fitting.pdf}
    \caption{\textbf{Parametric models fitted to real biosignals ground synthetic evaluation in physiological dynamics.} Neural parametric fits (dashed) overlaid on original recordings (black) for (\textbf{A})~EEG (CHB-MIT, channel FP1-F7, 10\,s segment), (\textbf{B})~ECG (MIT-BIH record 100, 10\,s segment, S-Q interval baseline), and (\textbf{C})~PPG/BVP (PPG-DaLiA subject S1, best-fit 5\,s segment). The parametric models capture the dominant oscillatory morphology of each signal type while providing closed-form ground truth for amplitude, frequency, and phase, enabling the fidelity evaluation framework used throughout this study.}
    \label{fig:fitting}
\end{figure}

\begin{table}[ht]
\caption{Mapping Physiological Phenomena to TimeSynth Evaluation Paradigms.}\label{tab:table1}
\begin{tabular}{@{}ll@{}}
\toprule
\textbf{Physiological Phenomenon} & \textbf{Synthetic Paradigm} \\
\midrule
Baseline wander (ECG, EEG) & Low-frequency drift (drift harmonic) \\
Respiration coupling (PPG) & Phase modulation \\
Heart rate variability (ECG) & Dual phase modulation + Markov switching \\
Neural oscillation variability (EEG) & Multi-frequency modulation + state switching \\
Sensor noise (all signals) & Additive noise (SNR 0--6 dB) \\
Exercise-induced changes & Non-stationary frequency shift \\
Arrhythmia onset (ECG) & Single state transition \\
Sleep-stage transitions (EEG) & State transition + Markov switching \\
Electrode artifacts (ECG, EEG) & Drift + noise \\
\bottomrule
\end{tabular}
\end{table}


\subsection*{Pointwise accuracy fails to reflect physiological fidelity}

Clinical decisions in cardiac monitoring, seizure detection, and respiratory assessment depend on the precise timing and rhythm of oscillatory events, not merely on waveform amplitude. If forecasting models inside digital twins achieve low prediction error while mistracking oscillatory phase or frequency, the resulting digital twin would appear reliable by standard benchmarks yet fail to capture the temporal dynamics that clinicians act on. Existing evaluation metrics for time series forecasting, including MSE, MAE, and their normalized variants, quantify only pointwise amplitude deviation and are therefore blind to whether a model preserves oscillatory timing or rhythm. To address this diagnostic gap, we introduce two additional fidelity metrics that capture the temporal properties existing metrics miss. Frequency error quantifies the mismatch in dominant oscillation rate between predicted and true signals, computed via peak detection with parabolic refinement on the power spectrum, with reliability filtering to exclude sequences lacking a clear periodicity (Methods; Supplementary Section~S6). Phase error quantifies temporal misalignment between predicted and true oscillatory cycles, computed from the instantaneous phase of the analytic signal with amplitude-based masking to exclude low-confidence regions near zero crossings (Methods; Supplementary Section~S6). Together with conventional amplitude error (MAE), these three diagnostics decompose forecast quality into the components that matter for clinical interpretation: a model may track amplitude accurately while drifting in phase, or preserve frequency content while distorting the oscillatory envelope, and only separate diagnostics can distinguish these failure modes.

We compared MAE against both phase error and frequency error across all 11 models and three signal families of increasing complexity (Fig.~\ref{fig:dissociation}a--f). Although MAE and phase error were strongly rank-correlated across models on all three signal families (Spearman $\rho = 0.93$ on drift-harmonic, $\rho = 0.90$ on single-phase modulation, and $\rho = 0.97$ on dual-phase modulation; $p < 0.001$ for each), a linear regression of phase error on MAE left a residual spread of $\mathrm{SD} = 9.7^{\circ}$, $3.5^{\circ}$, and $4.5^{\circ}$ respectively, with maximum residuals reaching $24.6^{\circ}$ on drift-harmonic signals (Supplementary Table~\ref{tab:phase_mae_residuals}). In other words, MAE ranks models in approximately the correct order, but cannot resolve the phase-fidelity differences that matter for physiological monitoring: models with comparable pointwise accuracy differed by up to $24^{\circ}$ in phase error on drift-harmonic signals and by more than $30^{\circ}$ within the degraded-performance cluster on dual-phase modulation signals (Fig.~\ref{fig:dissociation}a,c,e). Linear-family models consistently exhibited phase deviations exceeding $50^{\circ}$, whereas convolutional and patch-based architectures maintained phase error below $10^{\circ}$ (paired $t$-test, all pairwise differences Holm-corrected $p < 0.05$). A parallel dissociation emerged for frequency fidelity (Fig.~\ref{fig:dissociation}b,d,f): linear-family models showed frequency deviations up to 0.04~Hz on drift-harmonic signals and exceeding 0.05~Hz on dual-phase modulation signals at comparable MAE to architectures with near-zero frequency error. The frequency dissociation also increased with signal complexity, with transformer architectures showing the largest frequency deviations on dual-phase modulation signals (up to 0.20~Hz).

Had evaluation relied solely on MAE, the phase and frequency deficiencies of linear and transformer architectures would have remained undetected, and these models would have been classified as suitable for digital twin deployment. The subsequent analyses therefore employ all three fidelity diagnostics to systematically evaluate how architectural choice governs which dynamical properties a digital twin preserves under progressively challenging conditions.

\begin{figure}
    \centering
    \includegraphics[width=1.0\linewidth]{Figures/Fidelity/dissociation_combined.pdf}
    \caption{\textbf{Pointwise accuracy ranks models approximately but cannot resolve phase and frequency differences among models with comparable MAE.} Median MAE versus median absolute phase error (left column: \textbf{a}, \textbf{c}, \textbf{e}) and median MAE versus median absolute frequency error (right column: \textbf{b}, \textbf{d}, \textbf{f}) for all models across three signal families: drift-harmonic (\textbf{a,b}), single-phase modulation (\textbf{c,d}), and dual-phase modulation (\textbf{e,f}). Points are colored by architecture: CNN (orange), transformer (pink), MLP (green), and linear-family (gray). Although Spearman rank correlations between MAE and phase error are strong across all three families ($\rho = 0.90$--$0.97$, $p < 0.001$), residual phase-error spread at similar MAE levels reaches $24^{\circ}$ on drift-harmonic signals and $34^{\circ}$ within the degraded-performance cluster on dual-phase signals. CNN architectures consistently maintain phase error below $10^{\circ}$ and near-zero frequency error across all signal complexities. Linear-family models exhibit $>50^{\circ}$ phase error and elevated frequency deviation even at comparable MAE. Transformer architectures show the largest frequency deviations on dual-phase modulation signals. Both dissociations widen with signal complexity, demonstrating that conventional error metrics provide insufficient resolution of digital twin fidelity across multiple dynamical dimensions.}
    \label{fig:dissociation}
\end{figure}


\subsection*{Nonlinear architectures preserve oscillatory structure under clean signal conditions}

Having established that pointwise accuracy cannot resolve phase fidelity differences, we next asked which of the 11 architectural families preserve oscillatory structure when no external perturbation is present. This clean-signal evaluation serves as a performance ceiling: any architecture that fails to maintain amplitude, frequency, and phase fidelity under ideal conditions cannot be expected to do so under the noise, drift, and nonstationarity of real clinical signals.

On the most complex signals (dual-phase modulation), architectural families separated clearly by their ability to preserve dynamical structure (Fig.~\ref{fig:clean}a,b). Linear models showed phase errors often exceeding $50^{\circ}$ and elevated frequency errors, while convolutional and patch-based architectures maintained substantially lower error across all fidelity dimensions (paired $t$-test, Holm-corrected throughout this subsection). MICN\_Regre ($+58.60^{\circ}$, $p < 0.05$), MICN\_Mean ($+57.76^{\circ}$, $p < 0.05$), ModernTCN \cite{luo2024moderntcn} ($+52.73^{\circ}$, $p < 0.05$), and PatchTST ($+51.55^{\circ}$, $p < 0.05$) each improved phase accuracy by more than $50^{\circ}$ over the linear baseline, while Transformer ($-9.58^{\circ}$, $p < 0.05$) and Autoformer \cite{wu2021autoformer} ($-19.65^{\circ}$, $p < 0.05$) performed significantly worse than baseline (Fig.~\ref{fig:clean}b). FreMLP also exhibited a small but statistically significant deficit relative to Linear on dual-phase signals ($-1.61^{\circ}$, $p < 0.05$), consistent with spectral leakage under multi-frequency modulation. Frequency and amplitude distributions showed consistent architectural patterns (Supplementary Figs.~\ref{fig:supp_dual}--\ref{fig:supp_drift}). These results indicate that architectural inductive bias, rather than training objective alone, determines whether a model preserves the timing relationships embedded in physiological signals. Architectures that process the input through short, overlapping temporal windows, specifically patch-based and convolutional designs, inherently encode within-window periodicity that global attention and linear mappings average away.

The architectural ranking was not fixed, however, and shifted systematically with signal complexity. On drift-harmonic signals, where only a single slowly varying frequency is present, nearly all architectures including Transformer ($+45.38^{\circ}$, $p < 0.05$) surpassed the linear baseline for phase, while FITS ($-11.73^{\circ}$, $p < 0.05$) and Autoformer ($-26.14^{\circ}$, $p < 0.05$) degraded significantly; DLinear remained statistically indistinguishable from Linear at baseline ($-0.17^{\circ}$, Holm-corrected $p < 0.05$ but effect size negligible) (Fig.~\ref{fig:clean}c). This reversal indicates that global attention can preserve phase when the spectral structure is simple; it is specifically the multi-frequency modulation present in complex physiological signals that defeats it. NBeats exhibited the strongest signal-complexity dependence, rising from moderate improvement on dual-phase signals ($+14.22^{\circ}$, $p < 0.05$) to the best performer on single-phase ($+43.72^{\circ}$, $p < 0.05$) and drift-harmonic ($+53.50^{\circ}$, $p < 0.05$) signals. FreMLP showed a similar pattern, degrading below baseline on dual-phase ($-1.61^{\circ}$, $p < 0.05$) but improving substantially on drift-harmonic ($+50.09^{\circ}$, $p < 0.05$), consistent with spectral leakage under nonstationary modulation. Autoformer remained the weakest architecture across all signal families and fidelity dimensions ($p < 0.05$ for all pairwise comparisons versus Linear, Holm-corrected). Amplitude errors followed the same architectural hierarchy (Supplementary Figs.~\ref{fig:supp_dual}--\ref{fig:supp_drift}), with the largest absolute improvements on drift-harmonic signals where the linear baseline produced amplitude errors 2--3$\times$ those of the best-performing models.

In short, architectures that process the input through localized temporal windows (PatchTST, MICN, ModernTCN) consistently preserved oscillatory structure on complex signals, whereas architectures relying on global mappings or full-sequence attention preserved fidelity only when spectral structure was simple.

\begin{figure}
    \centering
    \includegraphics[clip, trim=0.3cm 6.5cm 0.3cm 2.5cm, width=1.0\textwidth]{Figures/Clean/clean.pdf}
    \caption{\textbf{PatchTST, MICN, and ModernTCN achieve $>50^{\circ}$ phase preservation on dual-phase signals, but architectural rankings reverse on simpler signals.} (\textbf{a})~Clean paradigm. Models observe history (blue) and predict future (green) against ground truth (orange). Temporal shifts in oscillatory peaks create phase misalignment invisible to pointwise metrics but critical for physiological monitoring. (\textbf{b})~Phase accuracy improvement ($|\Delta\text{phase}|$) relative to linear baseline for dual-phase modulation signals. MICN\_Regre ($+58.60^{\circ}$), MICN\_Mean ($+57.76^{\circ}$), ModernTCN ($+52.73^{\circ}$), and PatchTST ($+51.55^{\circ}$) achieve $>50^{\circ}$ improvement, while Transformer ($-9.58^{\circ}$) and Autoformer ($-19.65^{\circ}$) degrade below baseline, indicating that global attention disrupts phase on modulated signals. (\textbf{c})~Phase accuracy improvement for drift-harmonic signals. The ranking reverses: NBeats ($+53.50^{\circ}$) and ModernTCN ($+51.24^{\circ}$) lead, and Transformer ($+45.38^{\circ}$) now outperforms baseline, while FITS ($-11.73^{\circ}$) and Autoformer ($-26.14^{\circ}$) degrade. This reversal indicates that global attention preserves phase on simple spectral structure but fails under multi-frequency modulation. Error bars represent 95\% confidence intervals.}
    \label{fig:clean}
\end{figure}


\subsection*{Nonlinear models degrade gracefully under noise}

Real-world biosignals are invariably corrupted by sensor noise, such as motion artifacts in wearable PPG, electrode impedance fluctuations in ECG, and muscle contamination in EEG \cite{clifford2006advanced, sornmo2005bioelectrical}. A health digital twin deployed for ambulatory or bedside monitoring must maintain fidelity despite this corruption.

To assess noise robustness, all models were trained on clean signals and evaluated on test sequences with additive Gaussian noise at six signal-to-noise ratio (SNR) levels ranging from 40~dB (minimal corruption) to 1~dB (severe corruption), applied to the observed history while future targets remained clean (Methods).

On dual-phase modulation signals, PatchTST and MICN variants maintained the strongest phase advantage across noise conditions (Fig.~\ref{fig:noise}a,b). PatchTST led at low to moderate corruption ($+15.24^{\circ}$ at SNR~0, $p < 0.05$; $+8.22^{\circ}$ at SNR~4, $p < 0.05$), while MICN\_Regre ($+11.52^{\circ}$ at SNR~0, $p < 0.05$; $+2.44^{\circ}$ at SNR~6, $p < 0.05$) and MICN\_Mean ($+10.83^{\circ}$ at SNR~0, $p < 0.05$; $+1.84^{\circ}$ at SNR~6, $p < 0.05$) retained significant improvement across the full noise range. Under severe corruption (SNR~6), PatchTST's advantage narrowed to $+0.11^{\circ}$ (Holm-corrected $p > 0.05$, non-significant), but MICN variants remained significantly above baseline, demonstrating complementary robustness profiles within the same top tier. NBeats \cite{oreshkin2019n} retained modest but consistent positive improvement across noise levels (approximately $+1.5^{\circ}$ to $+1.9^{\circ}$, $p < 0.05$ at all levels). Autoformer showed the largest degradation across all noise levels ($-19.76^{\circ}$ to $-3.35^{\circ}$, $p < 0.05$ at all levels), FITS showed negative improvement throughout ($p < 0.05$), and DLinear exhibited a crossover from positive at low noise ($+1.82^{\circ}$ at SNR~0, $p < 0.05$) to negative under severe corruption ($-1.86^{\circ}$ at SNR~6, $p < 0.05$), confirming that linear-family architectures offer no robustness guarantee when sensor noise is high. All reported $p$-values are Holm-corrected for multiple comparisons.

The same architectural hierarchy held across signal families with informative shifts in relative ordering. On single-phase modulation signals, MICN\_Regre retained $+26.37^{\circ}$ at SNR~6 ($p < 0.05$) and PatchTST retained $+17.46^{\circ}$ ($p < 0.05$), both substantially ahead of all other architectures (Supplementary Fig.~\ref{fig:supp_noise_single}). Their relative ordering reversed, with MICN\_Regre outperforming PatchTST at high noise by $+8.91^{\circ}$, indicating that MICN's multi-scale convolution provides more stable retention under severe corruption on single-frequency signals. On drift-harmonic signals, the gap between architectures compressed: NBeats showed the best retention ($+23.17^{\circ}$ at SNR~6, $p < 0.05$), with PatchTST and MICN variants close behind (Supplementary Fig.~\ref{fig:supp_noise_drift}). Even on these simplest signals, Autoformer remained the worst performer across all three signal families ($p < 0.05$ at all noise levels). All cross-signal comparisons were Holm-corrected.

Across all signal families and noise conditions, PatchTST and MICN variants consistently emerged as the most reliable architectures for noise robustness, with PatchTST leading under low to moderate corruption and MICN providing more stable retention under severe corruption. A cardiac digital twin deployed in ambulatory settings would therefore need architecture selection matched to the expected noise regime, or an ensemble strategy that leverages the complementary strengths of PatchTST and MICN under different corruption levels.

\begin{figure}
    \centering
    \includegraphics[clip, trim=1.5cm 9.5cm 0.8cm 0.5cm, width=0.9\textwidth]{Figures/Noise/noisy.pdf}
    \caption{\textbf{PatchTST and MICN variants maintain phase accuracy under noise; linear models degrade systematically.} (\textbf{a}) Noisy paradigm. Models are trained on clean signals. At test time, Gaussian noise is added to observed history (blue) while the future remains clean. Predicted future (green) is compared against ground truth (orange). (\textbf{b}) Phase error difference across increasing noise levels (SNR 0 to 6, where higher number corresponds to more severe corruption) for dual-phase modulation signals. PatchTST leads at low to moderate noise ($+15.24^{\circ}$ at SNR~0). MICN variants retain significant improvement across all noise levels, including under severe corruption (SNR~6) where PatchTST's advantage narrows to non-significance. Autoformer shows consistent degradation ($-19.76^{\circ}$ to $-3.35^{\circ}$). Per-signal-family analyses are provided in Supplementary Figs.~\ref{fig:supp_noise_single}--\ref{fig:supp_noise_drift}. All differences marked with $*$ are Holm-corrected $p < 0.05$.}
    \label{fig:noise}
\end{figure}


\subsection*{Frequency shift adaptation is architecture-dependent}

Physiological rhythms are not stationary: heart rate increases during exercise, respiratory rate changes with posture and activity, and EEG frequency bands shift across arousal states \cite{clifford2006advanced, sornmo2005bioelectrical}. A digital twin must track these changes rather than assuming that the frequency characteristics observed during training will persist. If a forecasting model cannot adapt to shifted frequency distributions, the digital twin will continue generating predictions tuned to the old rhythm, missing the clinically informative transition.

To test generalization across frequency distributions, models trained on one frequency range were evaluated on shifted frequency distributions displaced by $-2$ to $+2$ range-widths from the training band, simulating clinically relevant transitions such as rest-to-exercise or waking-to-sleep (Methods).

On single-phase modulation signals, PatchTST and MICN variants showed the strongest phase advantages at the no-shift condition (Shift~0), with NBeats ($+43.72^{\circ}$, $p < 0.05$), PatchTST ($+43.23^{\circ}$, $p < 0.05$), MICN\_Mean ($+42.00^{\circ}$, $p < 0.05$), and MICN\_Regre ($+41.93^{\circ}$, $p < 0.05$) all exceeding $+40^{\circ}$ improvement over Linear (Fig.~\ref{fig:shift}a,b). These advantages declined under frequency shift but did not collapse uniformly. MICN\_Mean retained the most phase fidelity at Shift~$-2$ ($+9.35^{\circ}$, $p < 0.05$), followed by PatchTST ($+8.85^{\circ}$, $p < 0.05$), while NBeats dropped more steeply ($+5.56^{\circ}$, $p < 0.05$). At Shift~$+2$, all top models converged toward marginal improvement, with PatchTST at $+1.35^{\circ}$ ($p < 0.05$) and NBeats at $+0.67^{\circ}$ ($p < 0.05$). DLinear showed a modest but stable advantage across the $-2$ to $0$ range ($+1.76^{\circ}$ to $+2.43^{\circ}$, $p < 0.05$), consistent with its shift-invariant linear mapping. Transformer ($-9.01^{\circ}$ at Shift~0, $p < 0.05$) and Autoformer ($-15.71^{\circ}$, $p < 0.05$) degraded significantly even without any shift, indicating that global attention suppresses rather than tracks frequency-specific phase relationships. All $p$-values are Holm-corrected.

The severity of frequency shift degradation scaled with signal complexity. On dual-phase modulation signals, the collapse was more severe and roughly symmetric: MICN\_Regre dropped from $+58.60^{\circ}$ at Shift~0 ($p < 0.05$) to $-4.57^{\circ}$ at Shift~$-2$ ($p < 0.05$) and $-1.15^{\circ}$ at Shift~$+2$ ($p < 0.05$) (Supplementary Fig.~\ref{fig:supp_shift_dual}). On drift-harmonic signals, a directional asymmetry emerged: most architectures degraded more under negative shifts than positive ones, with MICN\_Mean falling to $-10.32^{\circ}$ at Shift~$-2$ ($p < 0.05$) but only $-1.94^{\circ}$ at Shift~$+2$ ($p < 0.05$) (Supplementary Fig.~\ref{fig:supp_shift_drift}), suggesting that downward frequency shifts pose a greater challenge than upward ones for models trained on higher-frequency bands. All $p$-values are Holm-corrected.

No architecture maintained more than approximately $10^{\circ}$ improvement at extreme shifts ($\pm 2$) on any signal family, revealing a fundamental boundary on how far current architectures can generalize beyond their training distribution. Within this boundary, PatchTST and MICN variants consistently provided the strongest adaptation at moderate shifts where linear-family and attention-based models had already degraded to baseline or below. A digital twin monitoring a patient whose heart rate shifts from resting (${\sim}60$~bpm) to moderate exercise (${\sim}100$~bpm), a shift well within the $\pm 1$ range, would maintain fidelity with PatchTST or MICN but lose phase tracking with linear or transformer-based architectures. Beyond moderate shifts, continuous model updating with incoming data will be essential for any architecture to maintain fidelity as physiological states evolve.

\begin{figure}
    \centering
    \includegraphics[clip, trim=1.5cm 11.0cm 0.8cm 0.5cm, width=0.8\textwidth]{Figures/Shift/shift.pdf}
    \caption{\textbf{PatchTST and MICN variants adapt to rhythm changes; linear and attention-based models fail across frequency shifts.} (\textbf{a}) Frequency distribution shift condition. Models trained on one frequency range (observed history, blue) forecast into a shifted frequency range (future, orange/green), simulating rest-to-exercise transitions or sleep-stage-dependent changes. (\textbf{b}) Phase error difference across frequency shift levels (Shift $-2$ to $+2$) for single-phase modulation signals. MICN variants, PatchTST, and NBeats achieve the largest improvements at Shift~0 ($+42^{\circ}$ to $+44^{\circ}$). Performance degrades at extreme shifts ($\pm 2$) but PatchTST and MICN\_Mean retain the most fidelity. Autoformer and Transformer degrade even at Shift~0 ($-15^{\circ}$ to $-9^{\circ}$). Per-signal-family analyses are provided in Supplementary Figs.~\ref{fig:supp_shift_dual}--\ref{fig:supp_shift_drift}. All differences marked with $*$ are Holm-corrected $p < 0.05$.}
    \label{fig:shift}
\end{figure}


\subsection*{PatchTST and ModernTCN adapt fastest to observable state transitions}

Among the most clinically consequential events a health digital twin must handle are abrupt physiological state changes: the onset of atrial fibrillation, the transition from interictal to ictal activity in epilepsy, or the shift between sleep stages \cite{sornmo2005bioelectrical}. For such events, the critical question is not whether a model can predict an unseen future transition, which no deterministic forecasting model can do without observable evidence, but how quickly a model adapts once evidence of the new state enters the observable history. A digital twin that requires 40 timesteps of post-transition data to recover phase tracking after a state change offers less clinical value than one that adapts within 6 timesteps.

To quantify health digital twin adaptability to state shifts, we introduced deterministic single-state frequency transitions at varying distances from the forecast boundary and measured how phase error evolved as the model received progressively more post-transition context (tags H2 to H40; Fig.~\ref{fig:transition}a,b). This design tests adaptation speed: when the transition falls within the observed history, models receive increasing amounts of new-state evidence, and we measure how much context each architecture requires to recover fidelity. When the transition falls in the unobserved future, no new-state evidence is available, and all models are expected to forecast based on the pre-transition state.

Architectures that process the input through short, overlapping temporal windows adapted fastest (Fig.~\ref{fig:transition}c). PatchTST and ModernTCN recovered to within $20^{\circ}$ of their no-transition baselines ($4.7^{\circ}$ and $5.8^{\circ}$, respectively) within 6 timesteps of post-transition context (Wilcoxon signed-rank test versus Linear, $p < 0.05$ at H6 for both, Holm-corrected). Their patch-based and convolutional designs encode local periodicity within short processing windows, allowing new temporal structure to overwrite learned representations rapidly once the transition enters the input. FreMLP and NBeats required approximately 15 timesteps to reach the same $20^{\circ}$ threshold ($p < 0.05$ at H15 for both), consistent with their broader temporal aggregation during inference. Transformer reached $15.4^{\circ}$ at H40 (significantly below Linear, $p < 0.05$), but required roughly 40 timesteps to fall below $20^{\circ}$, indicating that global attention distributes transition evidence across the full input context rather than concentrating it at the point of change. Linear-family models (Linear, DLinear, FITS, MLinear), Autoformer, and MICN never recovered below $20^{\circ}$, though this partly reflects their elevated no-transition baselines (Linear-family: $38^{\circ}$--$44^{\circ}$; Autoformer: $55^{\circ}$; MICN\_Mean: $89^{\circ}$). As expected, when transitions occurred in the unobserved future (tags F2 to F40), all models produced forecasts consistent with the pre-transition state, confirming that the paradigm tests adaptation to observable evidence rather than prediction of unseen events.

The adaptation speed hierarchy extended beyond phase to amplitude and frequency fidelity, though with dimension-specific differences (Supplementary Figs.~\ref{fig:supp_transition_mae}--\ref{fig:supp_transition_freq}). PatchTST and ModernTCN showed the fastest and most consistent recovery across all three fidelity dimensions ($p < 0.05$ versus Linear at H6 for both phase and frequency, Holm-corrected). MICN\_Mean exhibited a striking dissociation: strong frequency adaptation ($p < 0.05$ versus Linear at H10) but consistent amplitude degradation across all transition tags ($p < 0.05$ below Linear at H6 through H40), indicating that its multi-scale decomposition recovers rhythm tracking before magnitude preservation. This dimension-specific pattern illustrates why separate fidelity diagnostics are necessary for evaluating state transition performance.

For a digital twin receiving data at 10~Hz, PatchTST's 6-timestep adaptation corresponds to approximately 0.6 seconds of delayed state detection, while Transformer's 40-timestep requirement corresponds to 4 seconds. In a clinical setting where early detection of atrial fibrillation onset can trigger timely intervention, this difference in adaptation speed translates directly into the window of missed or delayed clinical response.

\begin{figure}
    \centering
    \includegraphics[clip, trim=1.5cm 7.5cm 0.8cm 1.0cm, width=1.0\textwidth]{Figures/State_Transition/state_transition.pdf}
    \caption{\textbf{PatchTST and ModernTCN adapt fastest to observable state transitions; adaptation speed varies by architecture.} (\textbf{a}) Single state transition from S1 to S2 occurring within the observed history (H-30, in-context), with the transition providing post-transition context before the forecast boundary. (\textbf{b}) Transition occurring in the unobserved future (F-40, no-context), with the forecast boundary preceding the state change. (\textbf{c}) Phase improvement over linear baseline as a function of post-transition context. Left region (H2 to H40, in-context): PatchTST and ModernTCN recover fastest, reaching strong improvement within 6 timesteps. FreMLP and NBeats require approximately 15 timesteps. Transformer adapts slowly, requiring roughly 40 timesteps. Linear-family models show minimal improvement regardless of available context. Autoformer and MICN\_Mean degrade below baseline. Right region (F2 to F40, no-context): all models produce forecasts consistent with the pre-transition state, as expected when no new-state evidence is available. Per-dimension analyses are provided in Supplementary Figs.~\ref{fig:supp_transition_mae}--\ref{fig:supp_transition_freq}. All differences marked with $*$ are Holm-corrected $p < 0.05$.}
    \label{fig:transition}
\end{figure}


\subsection*{PatchTST partially recovers stochastic switching dynamics where other architectures fail}

Cardiac rhythm alternates between stable and variable periods, sleep stages cycle with probabilistic dwell times, and autonomic tone fluctuates stochastically throughout the day \cite{clifford2006advanced, tong1990non}. The variability patterns that characterize conditions such as paroxysmal atrial fibrillation or cyclic alternating patterns in sleep are encoded in the switching statistics \cite{sornmo2005bioelectrical}. Unlike the deterministic transitions in the preceding analysis, stochastic switching provides within-window evidence of alternation patterns that a model could, in principle, learn to reproduce. To probe whether current forecasting architectures can recover these dynamics, we generated signals governed by a two-state Markov chain with transition probabilities $p \in \{0.10,\, 0.30,\, 0.50,\, 0.70,\, 0.90\}$ and assessed whether each model's forecast distribution matched the underlying switching statistics (Fig.~\ref{fig:markov}a).

Recovery was evaluated by fitting a two-state Gaussian HMM to windowed spectral features of the true history, decoding state sequences for both true and predicted futures, and comparing the resulting switching-probability distributions via symmetric KL divergence (Methods; Supplementary Section~S7.5). A model was classified as capturing the switching dynamics at a given $p$ if its symmetric KL divergence fell below 0.05 (Fig.~\ref{fig:markov}b).

PatchTST was the only architecture to recover switching dynamics at more than one transition probability (2/5), succeeding at $p = 0.70$ (KL $= 0.008$) and $p = 0.90$ (KL $= 0.046$). Both successes occurred at higher transition probabilities, where frequent switching provides more within-window evidence of alternation. Six architectures achieved recovery at a single probability each (1/5): DLinear (KL $= 0.022$ at $p = 0.90$) and Linear (KL $= 0.034$ at $p = 0.90$), FreMLP (KL $= 0.028$ at $p = 0.50$), ModernTCN (KL $= 0.014$ at $p = 0.70$), and MICN\_Mean (KL $= 0.016$ at $p = 0.50$) and MICN\_Regre (KL $= 0.026$ at $p = 0.50$). Five architectures, Autoformer (KL range: 0.521--2.020), FITS (0.202--0.574), MLinear (0.094--1.362), NBeats (0.061--0.324), and Transformer (0.102--1.242), failed across all settings (0/5). NBeats' complete failure is notable given its strong performance on the deterministic state-transition paradigm (recovering phase within 15 timesteps), highlighting that fast deterministic adaptation and probabilistic state recovery are distinct capabilities.

Stochastic switching was the most challenging paradigm tested: even PatchTST, which dominated the noise, shift, and deterministic state-transition paradigms, succeeded in fewer than half the conditions. The consistent failure across all four architectural families (linear, MLP, convolutional, and transformer-based) indicates that the limitation reflects a shared reliance on deterministic input-output mappings rather than a deficiency specific to any single design. The full KL divergence table and threshold sensitivity analysis are reported in Supplementary Table~\ref{tab:kl_full} and Supplementary Fig.~\ref{fig:kl_threshold_sensitivity}.

\begin{figure}[ht]
    \centering
    \includegraphics[clip, trim=2.6cm 11.5cm 0.8cm 1.0cm, width=0.9\textwidth]{Figures/markov/Markov.pdf}
    \caption{\textbf{PatchTST partially recovers switching dynamics; most architectures fail.} (\textbf{a}) Markov-switching signal alternating between states S1 and S2 with transition probability $p$. (\textbf{b}) Pass/fail summary across transition probabilities ($p = 0.10$ to $0.90$); blue denotes successful recovery (KL $< 0.05$), red denotes failure (KL $\geq 0.05$). PatchTST succeeds at $p = 0.70$ and $p = 0.90$ (2/5). Six architectures recover at a single probability each (1/5). Autoformer, FITS, MLinear, NBeats, and Transformer fail across all settings (0/5). Full KL values are reported in Supplementary Table~\ref{tab:kl_full}; threshold sensitivity analysis in Supplementary Fig.~\ref{fig:kl_threshold_sensitivity}.}
    \label{fig:markov}
\end{figure}


\subsection*{Localized temporal processing achieves the most balanced fidelity across paradigms}

The preceding analyses each isolate a single perturbation dimension. In practice, however, a digital twin deployed for patient monitoring faces noise, frequency shifts, and state changes simultaneously, and fidelity varies not only across individual dynamical properties but also across their interactions under combined perturbation \cite{clifford2006advanced, sornmo2005bioelectrical}. Strong performance on any single paradigm does not guarantee overall suitability. To identify architectures that maintain fidelity across all conditions, we performed a Pareto analysis, identifying models that cannot be improved on one dimension without sacrificing another, across all five paradigms: clean accuracy, noise robustness, shift robustness, state transition, and Markov fidelity (Fig.~\ref{fig:pareto}).

Each model was trained separately for each paradigm to isolate architectural effects from task-specific optimization. Four models reached the Pareto frontier (Fig.~\ref{fig:pareto}). MICN\_Regre and MICN\_Mean achieved frontier status through strong clean, noise, and shift performance but showed weakness in state transitions (normalized state-transition scores of 0.54 and 0.00, respectively; Supplementary Table~\ref{tab:pareto_scores}). ModernTCN reached the frontier through state-transition performance alone (0.97), trading off noise (0.31) and shift robustness (0.23). PatchTST was the only model maintaining performance above $0.8$ across all five dimensions (clean: 0.97, noise: 0.98, shift: 1.00, state transition: 1.00, Markov: 1.00), strictly dominating most alternatives. NBeats, despite being classified as dominated, achieved the third-highest mean score (0.85) with no dimension below 0.77 (Supplementary Fig.~\ref{fig:supp_radar_all}). No other architecture achieved above $0.8$ on more than three dimensions simultaneously. The full normalized score table and individual model radar profiles are provided in Supplementary Table~\ref{tab:pareto_scores} and Supplementary Fig.~\ref{fig:supp_radar_all}.

Across all paradigms, nonlinear architectures preserved fidelity more effectively than linear models but adapted only to observable changes. No current model reliably captured stochastic switching. Among all architectures tested, only PatchTST maintained balanced fidelity across noise corruption, distributional shifts, deterministic state transitions, and stochastic switching simultaneously, the combination of conditions characteristic of real patient monitoring. These results indicate that for digital twins operating under the multi-perturbation conditions of clinical deployment, architectures with localized temporal processing offer the most consistent fidelity, while the choice among them should be guided by which perturbation dimensions are most clinically relevant for the target application.

\begin{figure}
    \centering
    \includegraphics[width=0.8\linewidth]{Figures/Statistical/pareto_combined.pdf}
    \caption{\textbf{PatchTST is the only architecture maintaining balanced fidelity across all evaluation paradigms.} Radar plot comparing 11 models across five evaluation paradigms: clean accuracy, noise robustness, shift robustness, state transition, and Markov fidelity (0 to 1.0 scale). Solid polygons denote Pareto-optimal models; dashed gray polygons denote dominated models. PatchTST (blue) achieves balanced performance across all five dimensions ($>0.8$ on all axes), strictly dominating six models. MICN\_Regre (red) and MICN\_Mean (green) occupy the frontier but show weakness in state-transition adaptation. ModernTCN (purple) reaches frontier status through state-transition performance alone. Axes represent aggregate improvement over the linear baseline, averaged across all signals and conditions. Full normalized scores are reported in Supplementary Table~\ref{tab:pareto_scores}.}
    \label{fig:pareto}
\end{figure}
